\subsection{Independently certified upper bounds}

\begin{theorem}[Proved]\label{A:monoid-upper23}
There is an aperiodic monoid of order $23$ whose classifying space is
homotopy equivalent to $\CP^2$. The two supplied constructions numbered
$1412$ and $1651$ are isomorphic, and each has trivial automorphism group.
\end{theorem}
\begin{proof}
Let $I=\{0,1,2\}$ and $J=\{0,1,2,3,4\}$. We multiply transformation
pairs by
\[
 (\sigma,\tau)(\alpha,\beta)=(\sigma\alpha,\beta\tau).
\]
Take the identity, the fifteen constant pairs $x_{ij}$, and the seven
actors in the first pair of columns below. The second pair specifies the
other supplied construction, using its own labels $a,\ldots,g$.
\[
\begin{array}{c|cc|cc}
&\multicolumn{2}{c|}{1412}&\multicolumn{2}{c}{1651}\\
&\sigma&\tau&\sigma&\tau\\\hline
a&002&02200&010&04034\\
b&002&41144&020&04400\\
c&010&04034&020&21122\\
d&202&03300&020&30033\\
e&202&12211&022&01100\\
f&202&40044&022&32233\\
g&212&00334&212&00334
\end{array}
\]
The indicated sets are closed under this associative multiplication. Our
fresh constructor and separate verifier check all $529$ products and
$12\,167$ associativity triples for each table. The literal tables use
labels $0=1$, $x_{ij}=1+5i+j$, and $a,\ldots,g=16,\ldots,22$.
They are in \path{monoid/certificates/cp23_candidate_1412.json} and
\path{monoid/certificates/cp23_candidate_1651.json}.

We give the homotopy certificate for the first table. Put $R=\Z M$ and
define projective left $R$-modules
\[
 P=Rx_{00},\quad E=R(c-x_{00}),\quad F=R(1-c),\quad G=R(a-x_{00}).
\]
Each displayed generator is a ring idempotent. The exact resolution is
\[
0\longrightarrow P\xrightarrow{d_4}G\oplus G
 \xrightarrow{d_3}P\oplus F\xrightarrow{d_2}E
 \xrightarrow{d_1}P\xrightarrow{\varepsilon}\Z\longrightarrow0,
\]
where the $R$-linear maps are determined by
\begin{align*}
d_1(c-x_{00})&=x_{10}-x_{00},\\
d_2(x_{00},0)&=x_{03}-x_{00},\\
d_2(0,1-c)&=c+d-g-x_{00},\\
d_3(a-x_{00},0)&=(x_{00}-x_{20},b+e-f-x_{01}),\\
d_3(0,a-x_{00})&=(0,a-x_{00}),\\
d_4(x_{00})&=(x_{02}-x_{00},-x_{02}+x_{00}).
\end{align*}
For $P$ use the basis $x_{00},x_{10},x_{20}$. Bases for the other modules
are $z-zx_{00}$ for $E,G$ and $z-zc$ for $F$, where respectively
\begin{align*}
E:\quad&z\in\{x_{03},x_{04},x_{13},x_{14},x_{23},x_{24},c,d,f,g\},\\
F:\quad&z\in\{1,x_{01},x_{02},x_{11},x_{12},x_{21},x_{22},a,b,e\},\\
G:\quad&z\in\{x_{02},x_{12},x_{22},a\}.
\end{align*}
The verifier reconstructs these bases inside $\Z^{23}$ and checks every
translate of each generating idempotent has integer coordinates in its
claimed basis. Thus these are the whole principal projectives. It checks
all module actions, all differential equivariance equations, and all zero
composites. The abelian ranks of the terms are $(3,10,13,8,3)$; the ranks
of $\varepsilon,d_1,d_2,d_3,d_4$ are $(1,2,8,5,3)$; every nonzero Smith
invariant is $1$. Each differential image is consequently saturated,
contained in the next kernel, and of the same rank. It equals that kernel.
The last differential is injective. This proves exactness in every degree.

For a ring idempotent $e$, $\Z\otimes_R Re$ is the image of multiplication
by $\varepsilon(e)$ on $\Z$. Consequently the coinvariant terms are
$\Z,0,\Z,0,\Z$, with all differentials zero. We obtain integral homology
$\Z$ in degrees $0,2,4$, and zero in every other degree.

Extend the matrix with rows $(1,1),(0,1),(-1,0),(1,0)$ by zero row and
column at index zero, and call it $L$. The normalized cochain
\[
u(m,n)=L\bigl(\tau_m(0),\sigma_n(0)\bigr)
\]
is an integral two-cocycle, as checked on every triple. We compare the
projective resolution with the normalized left bar resolution. Its
contracting homotopy on a term $z[m_1|\cdots|m_d]$ inserts $z$ at the
front with coefficient $1$ if $z\ne1$, and gives zero if $z=1$.
Starting from $F_0(x_{00})=x_{00}[\,]$, we set
\[
 F_d(e)=e\,s(F_{d-1}d_d(e))
\]
on each projective generator and extend $R$-linearly. Idempotent support
makes the formula well defined. The separate verifier checks every chain
equation and every equivariance equation of the stored comparison.
The comparison has $(3,40,172,448,324)$ sparse terms in degrees $0$--$4$.
Integral evaluation on the coinvariant generators gives
\[
\langle u,F_2(x_{00},0)\rangle=-1,
\qquad \langle u\smile u,F_4(x_{00})\rangle=1.
\]
Thus the degree-two class and its square are primitive.

Every band element is idempotent, hence trivial in the universal group;
the relation $mx_{00}\in K$ then makes every element of $M$ trivial there.
Thus $BM$ is simply connected. We recall the recognition argument, also
covered by \cite[Proposition~4C.3]{CP2-Hatcher2002}. A Hurewicz generator
$S^2\to BM$ has a $3$-connected mapping-cylinder pair: relative
Hurewicz in degree three applies after the degree-two isomorphism and
the vanishing of relative $H_3$. Relative Hurewicz in degree four then
represents the sole relative homology generator by a map of a four-cell.
This gives an integral homology equivalence
$S^2\cup_h e^4\to BM$ between simply connected CW complexes, hence a
homotopy equivalence. The primitive square makes the Hopf invariant of
$h$ equal to $\pm1$. Since $\pi_3(S^2)$ is generated by a Hopf map,
this cofiber is $\CP^2$. The verifier separately checks stabilization of
every element's powers, proving aperiodicity.

The row permutation $(0\ 2)$ and column permutation $(1\ 2)(3\ 4)$
conjugate the first actor set to the second. Their element permutation is
\[
\begin{split}
(&0,11,13,12,15,14,6,8,7,10,9,1,3,2,5,4,\\
 &20,21,22,17,18,19,16).
\end{split}
\]
Our independent isomorphism program checks all $720$ row/column
permutations for every ordered pair of the two models. An isomorphism
must preserve the characteristic minimal band; its restriction there is
a row and column permutation, and faithfulness of the actor representation
determines its values elsewhere. This search is therefore exhaustive. It
finds exactly one isomorphism for each ordered pair, including each
automorphism pair. Neither listed monoid is anti-isomorphic to either
listed monoid, because opposition swaps the unequal minimal-band row and
column counts. Their opposites belong to the same class when opposition
is allowed.
\end{proof}

\begin{lemma}[Proved: constant minimal-band compression]
\label{A:band-action-compression}
Let $M$ have a proper rectangular-band minimal ideal $K=I\times J$,
and let $S$ be a nonempty finite right $M$-set on which every $k\in K$
acts constantly, with value $c_k$. Let $C$ have objects $a,b$, corners
$M,\{1_b\}$, cross hom-set $C(a,b)=S$, and $C(b,a)=\varnothing$.
There is a monoid
\[
 N=(M\setminus K)\amalg(I\times S),\qquad BC\simeq BN.
\]
\end{lemma}
\begin{proof}
Let $\rho:K\to I$ be the row map and let $\sigma_m:I\to I$ be the
left row action. Set $\eta(k)=(\rho(k),c_k)$. Preserve old products
outside $K$, replace old products in $K$ by their $\eta$-images, and set
\[
m(i,s)=(\sigma_m(i),s),\quad (i,s)m=(i,sm),\quad
(i,s)(j,t)=(i,t).
\]
The action equations give
$\eta(mk)=(\sigma_m\rho(k),c_k)$ and
$\eta(km)=(\rho(k),c_km)$; these verify all associativity cases
meeting the new band. The other cases are inherited from $M$. Thus
$q:M\to N$, equal to the identity outside $K$ and to $\eta$ on $K$,
is a unital homomorphism.

Fix $e=(i_0,s_0)$ and let $D$ be the full subcategory of
$\operatorname{Kar}(N)$ on $1,e$. The object $e$ retracts from $1$.
The inclusion $N\to D$ induces a homotopy equivalence by
Quillen's Theorem A \cite[\S1, Theorem~A]{CP2-Quillen1973}.
Indeed, its comma category over $e$ has objects $x\in eN$ and
morphisms $m:x\to y$ satisfying $ym=x$. The components
$x:x\to e$ define a natural transformation from the identity to
the constant functor at $e$: one has $ex=x$, and naturality is
exactly $ym=x$. The comma category over $1$ has a terminal object.
Thus both comma categories are contractible. Define $F:C\to D$ by
$a\mapsto1$, $b\mapsto e$, $m\mapsto q(m)$ and
$s\mapsto(i_0,s)$. The displayed product rules verify its functor laws.

The comma category $(F\downarrow e)$ has terminal object $(b,e)$.
In $(F\downarrow1)$ the objects over $a$ are $n\in N$, with arrows
$n\to n'$ labelled $m\in M$ when $n'q(m)=n$; the objects over $b$
are one sink for each $i\in I$. There is one arrow $(i,s)$ to sink $i$.
For a right $M$-set $X$, write $A(X)$ for the category of arrows
$x\to y$ labelled $m$ with $ym=x$. Its degree-$d$ nerve is naturally
$X\times M^d$, so it preserves colimits and injections. As a right
$M$-set, $N=M\amalg_K(I\times S)$. Each row category is a sieve,
so the comma nerve is obtained by attaching a cone to each row nerve.
Its realization is the pushout
\[
 BA(M)\amalg_{BA(K)}\coprod_{i\in I}\operatorname{Cone}(BA(S)).
\]
The map $BA(K)\to BA(M)$ is a simplicial subcomplex inclusion, hence
a cofibration, so this is a homotopy pushout. The category $A(M)$ has
terminal object $1$. Each component $A(e_iM)$ of $A(K)$ contracts by
the natural transformation $x\to e_i$ labelled $x$: naturality is
exactly $ym=x$. Therefore $BA(K)\to I$ is a homotopy equivalence.
Contracting the cones now shows that the comma category is contractible.
Quillen's theorem gives $BC\simeq BD\simeq BN$.
\end{proof}

\begin{theorem}[Proved]\label{A:category-upper22}
There is a two-object category with $22$ arrows, identities included,
and classifying space $\CP^2$. Both supplied $22$-arrow constructions
are isomorphic and have been independently checked.
\end{theorem}
\begin{proof}
In either monoid $N$ from Theorem~\ref{A:monoid-upper23}, retain the
eight elements outside the band and the ten band elements in rows $0,2$.
This is an $18$-element submonoid $P$: every outside actor preserves
$\{0,2\}$, and a product of two outside actors cannot be the constant
row map with value $1$. Indeed an outside actor taking the value $1$
has singleton inverse image of $1$, so a composite of two such
nonconstant maps cannot be constantly $1$.

Set $M=P^{\mathrm{op}}$ and let $S=I$ act on the right by the original
row transformations. Every element of its $5\times2$ minimal band
acts constantly. The resulting act-cone category has $18+1+3=22$
arrows. Compression in Lemma~\ref{A:band-action-compression} replaces
the $5\times2$ band by a $5\times3$ band and recovers exactly
$N^{\mathrm{op}}$. Its four product types agree with the specified
multiplication. Thus $BC\simeq B(N^{\mathrm{op}})\simeq\CP^2$.

The files \path{monoid/certificates/cp22_candidate_1412.json} and
\path{monoid/certificates/cp22_candidate_1651.json} give every source,
target, identity and composition, as well as the full comparison target
and functor. Undefined composites are recorded as $-1$. The independent
checker verifies $972$ action equations, $382$ composable-pair functor
equations and $6862$ composable associativity triples for each category.
The isomorphism of the source monoids preserves the retained row set;
its restriction on the corner and its row permutation on the cross
arrows give an isomorphism of the two categories, checked on every
composition.

The actual comparison functor has target the two-object category
$\operatorname{Kar}(N^{\mathrm{op}})\{1,x_{00}\}$, with $32$ arrows.
This corrects the prompt's wording ``onto the $23$-element monoid.''
The comparison is the proved span of homotopy equivalences; the upper
bound is unaffected. Neither theorem in this subsection asserts a
matching lower bound or a global classification of all minimal models.
\end{proof}
